%--------------------------------------------------------------------%
%
% Bab IV — Evaluasi dan Pembahasan
%
% @author Petra Novandi
%
%--------------------------------------------------------------------%
%
% Berisi tujuan pengujian, skenario, hasil pengujian, dan pembahasan.
% Ganti teks placeholder dengan isi sebenarnya.
%
%--------------------------------------------------------------------%

\chapter{Evaluasi dan Pembahasan}

Bab ini mempresentasikan pengujian solusi yang dirancang pada Bab
III dan membahas apakah tujuan serta rumusan masalah telah
terpenuhi. Hasil yang ditampilkan harus dapat dipertanggungjawabkan
secara objektif.

\section{Tujuan Pengujian}

Tuliskan tujuan pengujian yang spesifik dan terukur. Setiap tujuan
sebaiknya berkaitan langsung dengan tujuan tugas akhir dan
rumusan masalah di Bab I.

\blindtext

\section{Skenario Pengujian}

Jelaskan metode, lingkungan, data uji, dan langkah-langkah
pengujian yang dilakukan. Skenario harus cukup jelas agar pengujian
dapat direplikasi oleh pembaca.

\blindtext

\section{Hasil Pengujian}

Sajikan hasil pengujian secara faktual menggunakan tabel, grafik,
atau angka. Hindari interpretasi panjang di subbab ini; fokuskan
pada data dan temuan yang diperoleh.

    \subsection{Subbab}

    \blindtext

    \subsubsection{Subsubbab}

    \blindtext

\blindtext

\section{Pembahasan}

Interpretasikan hasil pengujian dan kaitkan dengan rumusan masalah
serta tujuan tugas akhir. Jelaskan apakah solusi yang diusulkan
berhasil, keterbatasannya, dan perbandingannya dengan studi terkait
jika relevan.

\blindtext
